\section{Future Work}
\label{sec:future-work}

Each item below is scoped precisely enough to seed a follow-up run: the
experiment, the hypothesis it tests, how it builds on this round, and the
resources required. The reference compute profile is the one this run used
(a shared Slurm cluster's free partition of NVIDIA L40 48\,GB GPUs, jobs
under 30 minutes); items exceeding it say so explicitly.

\begin{enumerate}
  \item \textbf{Full-scale, non-subsampled Numerai verdict (removes the
    regime-drift scope limit).} Re-run the exp-004 protocol on the full
    training data and feature set at $B \approx 4096$, with F1/F5/F9
    machinery re-registered on fresh eras. \emph{Hypothesis:} the negative
    verdict persists when the baseline edge returns to the practitioner
    sanity band---i.e., the result is not an artifact of the low-edge
    subsampled regime (limitation~1). \emph{Builds on:} the frozen
    protocol, thresholds-as-code, and the audit's independent-shard
    tooling. \emph{Resources:} est.\ 10--20 L40-hours, or A100-class
    GPUs for comfortable $B{=}4096$ per-sample-gradient memory;
    \textbf{exceeds} the free-partition 30-minute cap---needs a paid or
    time-unconstrained allocation (order \$50--100 of cloud A100 time).
  \item \textbf{Spectral operating-point sweep at a declared enlarged
    matched budget (closes the F12 asymmetry).} Give \emph{both} arms a
    12-trial budget (spectral: mp\_factor $\times$ mode $\times$ $B$;
    AdamW re-swept at the same enlarged compute), pre-registered as a new
    protocol so the matched-compute design is redefined rather than
    broken. \emph{Hypothesis:} no spectral configuration in a reasonable
    neighborhood of the prior-project priors recovers the baseline---
    upgrading ``no evidence of benefit under the affordable tuning
    budget'' toward ``no benefit at any tested operating point''.
    \emph{Builds on:} exp-003's measured per-trial costs (63\,s/spectral
    run) and exp-001's registered diagnostics. \emph{Resources:}
    $\sim$30 GPU-minutes on one L40 (fits the free partition as two
    jobs); zero new data.
  \item \textbf{True second setting: OHLCV dataset, both architectures
    (converts ``architecture consistency'' into ``second setting'').}
    Build the pre-specified OHLCV protocol (public daily equities/crypto,
    purged temporal split, $\geq$90\% coverage after cleaning) and run
    MLP and GRU arms with the identical verdict machinery.
    \emph{Hypothesis:} the class-level harm of agreement-style filtering
    generalizes beyond Numerai's obfuscated within-era construct to a
    natural sequential financial task. \emph{Builds on:} exp-002's
    validated vmap-GRU stack and exp-005's within-era machinery.
    \emph{Resources:} $\sim$1--2 L40-hours total across 2--4 jobs (fits
    the free partition); data is free to download; $\sim$2--4 hours of
    engineering for the split/cleaning protocol.
  \item \textbf{Second loss family with target-sensitive engagement
    diagnostics (tests the theorem's boundary).} Repeat the comparison
    with a corr-aligned per-era loss and/or a multi-output head, where
    Theorem~\ref{thm:target-independence} does not apply (the audit's
    scoping controls confirm 2-output MSE breaks the invariance), using
    new engagement diagnostics (unnormalized-gradient spectra,
    residual-magnitude-weighted similarity, sign-pattern statistics of
    $D$). \emph{Hypothesis:} where the engagement spectrum is
    target-sensitive, the filter's behavior (and possibly its effect)
    changes---directly testing whether target-blind selection is the
    operative harm mechanism. \emph{Builds on:} the theorem and the C1
    confound calibration. \emph{Resources:} $\sim$30--60 L40-minutes
    across 2--3 jobs (fits the free partition).
  \item \textbf{Muon and GBT context baselines (the designated cuts).}
    Add a tuned Muon arm at matched budget
    \citep{gorishniy2026benchmarkingoptimizers} and a LightGBM-class GBT
    reference. \emph{Hypothesis:} the tuned-AdamW baseline is not the
    weak link---the filter's deficit persists against the current tabular
    bar. \emph{Builds on:} the frozen t07 pipeline; pure arm additions.
    \emph{Resources:} $\sim$30--60 L40-minutes ($\approx$2 experiment
    slots); no new data.
  \item \textbf{MP-null theory under cross-sectional correlation.}
    Simulation (and, where tractable, analytical) characterization of the
    Marchenko--Pastur bulk edge for per-sample gradients with financial
    factor structure, beyond the C2 permutation control---the shifted-edge
    question raised at design review. \emph{Hypothesis:} the hard-mode
    threshold systematically misclassifies factor-driven directions as
    ``signal'' under realistic cross-sectional correlation, formalizing
    the C1(c) confound result. \emph{Builds on:} the saved spectra
    (\texttt{spectra.npz}) and C1 test harness. \emph{Resources:} CPU
    only---days of workstation CPU, no GPU, no cluster constraint.
  \item \textbf{Small closure items} (bundle into one $\sim$25-minute L40
    job in any follow-up): (a) the $\leq$4-trial GRU tuning sweep and
    exp-005 re-run ($\sim$10 GPU-minutes; limitation~3); (b)
    feature-neutral correlation reporting via saved per-row predictions
    ($\sim$6 GPU-minutes); (c) a higher-power C5 era-identity probe
    ($\sim$5 GPU-minutes); (d) extra seed pairs to tighten the C4
    equivalence bound below $\pm 0.002$ ($\sim$80\,s per seed pair); and
    (e) the free held-out-seed re-bucketing of the tail-era split
    (CPU-only re-analysis of existing CSVs).
\end{enumerate}
